% TOP_PROBLEM: {"id": 3103, "title": "Reconstruction conjecture", "category_id": 3, "category": {"id": 3, "name": "graph_theory", "display_name": "Graph Theory", "description": "Problems involving graphs, networks, and their properties.", "slug": "graph-theory", "order_index": 3, "created_at": "2026-07-31T15:26:25.671Z"}, "status": "open", "problem_number": "OPG-658", "set_id": 12}
% FIRST_POSTED: 2026-10-01
% ATTEMPT_STATUS: unresolved
% COMPLETION: 5
% MODEL: GPT 6 Astra Ultra
% UnsolvedMath ID 3103; source catalogue code OPG-658
% !TeX program = lualatex
% Research attempt, 2026-10-01. Canonical statement preserved verbatim.
\documentclass[11pt,a4paper]{article}
\usepackage[margin=25mm]{geometry}
\usepackage{fontspec}
\setmainfont{lmroman10-regular.otf}[BoldFont=lmroman10-bold.otf,
 ItalicFont=lmroman10-italic.otf,BoldItalicFont=lmroman10-bolditalic.otf]
\usepackage{amsmath,amssymb,mathtools}
\usepackage{unicode-math}
\setmathfont{latinmodern-math.otf}
\AtBeginDocument{\let\setminus\smallsetminus}
\usepackage{xurl,hyperref}
\hypersetup{colorlinks=true,urlcolor=blue}
\setcounter{secnumdepth}{0}
\setlength{\parindent}{0pt}
\setlength{\parskip}{5pt}
\setlength{\emergencystretch}{3em}
\begin{document}
\section*{Reconstruction conjecture}\label{3103}
{\small UnsolvedMath ID 3103 · OPG-658 · Graph Theory · OpenGarden}
\subsection{Definitions and mathematical statement}
Let $G=(V,E)$ be a finite simple undirected graph. For $v\in V$, the
vertex-deleted graph $G-v$ is the induced graph on $V\setminus\{v\}$:
the vertex and all its incident edges are removed. Write $[H]$ for the
isomorphism class of an unlabelled graph $H$, where graph isomorphism is
a vertex bijection preserving adjacency in both directions.

The vertex deck of $G$ is the multiset
\[
 \mathcal D(G)=\{\!\{[G-v]:v\in V\}\!\}.
\]
There is one card for each vertex, and repeated isomorphism classes are
retained with their multiplicities. Neither the deleted vertex's identity
nor a consistent labelling across the cards is supplied.

The reconstruction conjecture states that for any finite simple graphs
$G,H$ with at least three vertices,
\[
 \mathcal D(G)=\mathcal D(H)\quad\Longrightarrow\quad G\cong H.
\]
Equality of decks includes equality of multiplicities, hence their
cardinalities and the orders of the original graphs. Both connected and
disconnected graphs, including isolated vertices, are allowed.

The claim is uniqueness up to isomorphism, not recovery of an original
vertex labelling. A counterexample would be a pair of nonisomorphic
simple graphs on the same number of at least three vertices with identical
decks. Matching a selection of cards, or merely matching counts of small
subgraphs, is weaker than deck equality. Directed graphs, parallel edges
and loops are outside this formulation, since altering the graph class
changes the reconstruction question.
\subsection{Short English statement}
Is every finite simple graph with at least three vertices uniquely determined up to isomorphism by its multiset of vertex-deleted graphs?
\subsection{Research attempt}
\paragraph{Scope and literature check.}
Kelly's original paper [A1] contains the subgraph-counting observation behind
reconstruction arguments. We prove its elementary counting mechanism and use
it to give a complete reconstruction algorithm for regular graphs. This is
an established special case, not a new resolution. The source page [S2] and
Kelly's paper were checked on 1 October 2026. We then isolate the loss of
vertex alignment that prevents this algorithm from extending automatically
to graphs with varying degrees.

\paragraph{What an unlabelled deck actually determines.}
Let the deck have $n\ge3$ cards $D_1,\ldots,D_n$, with $m_i$ edges in card
$i$. Every original edge occurs in exactly $n-2$ cards, so
\[
                  m=|E(G)|=\frac{\sum_i m_i}{n-2}.
\]
The deleted vertex corresponding to card $i$ had degree $m-m_i$.
Consequently the multiset of vertex degrees is reconstructible even though
no consistent vertex labels across cards are supplied.
More generally, for a fixed graph $F$ with $k<n$ vertices, let $i(F,J)$ count
vertex subsets of $J$ inducing a copy of $F$. Every such subset in $G$
survives in exactly $n-k$ cards, hence
\[
                 i(F,G)=\frac{\sum_i i(F,D_i)}{n-k}.
\]
This identity includes multiplicities and is valid for disconnected $F$ too.
It is a counting identity, not a way to identify which occurrences in two
cards represent the same original vertices.

\paragraph{A complete regular-graph reconstruction.}
Suppose the degree multiset just computed is constant, equal to $r$. Then
\emph{every} graph with this deck is $r$-regular, so regularity is recognized
rather than being assumed of a competing realization. Choose any one card
$D=G-v$. Its vertices have degrees either $r$ or $r-1$: deleting $v$ lowers
precisely its neighbours' degrees by one. Adjoin a new vertex and connect it
to every degree-$(r-1)$ vertex of $D$, and to no degree-$r$ vertex. The result
is isomorphic to $G$. There is no arbitrary matching step and no hidden
choice among same-degree vertices. Any isomorphism between two copies of the
card preserves their degrees and therefore extends to the reconstructed graph.
For $r=0$ the instruction joins the new vertex to none; for $r=n-1$ it joins
to every card vertex. Thus empty and complete graphs are covered as well.

\paragraph{Why a degree-only generalization is insufficient.}
For a general graph, a card vertex of degree $d$ could have had original
degree $d$ and been a non-neighbour of the missing vertex, or original degree
$d+1$ and been a neighbour. The deck determines how many vertices have each
original degree but does not assign those degrees to the vertices of this
particular card. Choosing such assignments independently across cards is not
justified. As a basic stress test, $C_6$ and the disjoint union of two triangles
both have six vertices, six edges, and constant degree two; those numerical
invariants alone cannot distinguish them. Their actual cards do distinguish
them: a card of $C_6$ is a five-vertex path, whereas a card of the other graph
is a triangle plus an edge. The regular reconstruction above uses this full
card structure, not merely the degree multiset.

\paragraph{Verification and first remaining gap.}
A finite enumeration checks the reconstruction procedure on every labelled
regular graph with $3\le n\le5$ and every possible deleted vertex. Degrees
recovered from deck edge counts agree with the originals, and adding the
missing vertex by the deficit rule restores the original labelled adjacency
on surviving vertices. The proof handles all orders independently of that
check. The factor $n-2$ explains the lower bound on $n$; it is not legitimate
to use this formula on two vertices. For nonregular graphs, the unresolved
step is a consistent assignment of original degrees and neighbourhoods that
is forced by the entire unlabelled deck. No uniqueness theorem for that step
is proved here.

\paragraph{Reproducibility.}
The finite calculations above are implemented in
\texttt{scripts/check\_attempt\_batch\_2026\_10\_01.py}; run it with Python 3
from the repository. It uses exact integer arithmetic and no external packages.

\paragraph{Final status.}
\textbf{UNRESOLVED.} Regular graphs are reconstructed with an explicit,
verified procedure and the obstruction to the naive nonregular extension is
identified. The universal reconstruction conjecture is not settled.
Completion estimate: 5\%.

\subsection{Sources}
\begin{itemize}
\item[S1] ULAM.AI and the UnsolvedMath Contributors, supplied problems.json, record 3103 (OPG-658), OpenGarden collection. Catalogue identity and source transcription; CC BY 4.0.
\url{https://huggingface.co/datasets/ulamai/UnsolvedMath}.
\item[S2] Open Problem Garden, Reconstruction conjecture. Original problem and discussion; the mathematical formulation below is expanded from the supplied source record.
\url{http://www.openproblemgarden.org/op/reconstruction_conjecture}.
\item[A1] P. J. Kelly, \emph{A congruence theorem for trees}, Pacific Journal of Mathematics 7 (1957), 961--968, especially the initial counting lemmas. \url{https://msp.org/pjm/1957/7-1/pjm-v7-n1-p14-p.pdf}.
\end{itemize}
\end{document}
